\section*{अध्याय \ref{ch:b2:aromatic-substitution} --- ऐरोमैटिकता और इलेक्ट्रॉनरागी ऐरोमैटिक प्रतिस्थापन}

\begin{solution}{exo:b2:aromatic-substitution:1}
\ce{HNO3 + 2H2SO4 -> NO2+ + H3O+ + 2HSO4-}: नाइट्रिक अम्ल प्रोटॉनित होता है, फिर जल खो देता है।
\end{solution}

\begin{solution}{exo:b2:aromatic-substitution:2}
टॉलूईन 2-ब्रोमोटॉलूईन और 4-ब्रोमोटॉलूईन देती है: मेथिल ऑर्थो और
पैरा स्थितियों की ओर निर्देशित करता है। नाइट्रोबेन्ज़ीन 3-ब्रोमोनाइट्रोबेन्ज़ीन देती है, अधिक धीमे: नाइट्रो समूह
मेटा निर्देशित करता है।
\end{solution}

\begin{solution}{exo:b2:aromatic-substitution:3}
\ce{-CH3}: सक्रियक, ऑ/पै। \ce{-OH}: प्रबल सक्रियक, ऑ/पै। \ce{-Cl}: निष्क्रियक, ऑ/पै।
\ce{-NO2}: प्रबल निष्क्रियक, मेटा। \ce{-COCH3}: निष्क्रियक, मेटा। \ce{-NHCOCH3}:
सक्रियक (मध्यम), ऑ/पै।
\end{solution}

\begin{solution}{exo:b2:aromatic-substitution:4}
ऐसीटोफ़ीनोन (1-फ़ेनिलएथेन-1-ओन), \ce{C6H5COCH3}, \omterm{def:b2:aromatic-substitution:friedel-crafts}{ऐसिलियम आयन} \ce{CH3CO+} के माध्यम से।
\end{solution}

\begin{solution}{exo:b2:aromatic-substitution:5}
प्राथमिक कार्बधनायन (या उसका संकुल) हाइड्राइड विस्थापन द्वारा द्वितीयक
आइसोप्रोपिल धनायन में पुनर्विन्यस्त होता है, जो वलय का ऐल्किलन करता है। प्रोपिलबेन्ज़ीन: प्रोपेनॉयल क्लोराइड से
\omterm{def:b2:aromatic-substitution:friedel-crafts}{फ़्रीडेल--क्राफ़्ट्स ऐसिलन} (कोई पुनर्विन्यास नहीं), फिर कीटोन \ce{C=O} का \ce{CH2} में अपचयन।
\end{solution}

\begin{solution}{exo:b2:aromatic-substitution:6}
\ce{-OH} समूह एक प्रबल मेसोमेरी दाता है: वलय इतना सक्रिय है कि आण्विक ब्रोमीन
पर्याप्त इलेक्ट्रॉनरागी है, और तीनों ऑर्थो और पैरा स्थितियाँ प्रतिस्थापित हो जाती हैं।
\end{solution}

\begin{solution}{exo:b2:aromatic-substitution:7}
3-ब्रोमोनाइट्रोबेन्ज़ीन: पहले नाइट्रीकरण, फिर ब्रोमीनन (नाइट्रो मेटा निर्देशित करता है)। 4-ब्रोमोनाइट्रोबेन्ज़ीन:
पहले ब्रोमीनन, फिर नाइट्रीकरण (ब्रोमीन ऑर्थो/पैरा निर्देशित करता है), और पैरा समावयवी को
ऑर्थो से पृथक कीजिए।
\end{solution}

\begin{solution}{exo:b2:aromatic-substitution:8}
नाइट्रीकारी मिश्रण अत्यंत इलेक्ट्रॉन-समृद्ध ऐनिलीन का ऑक्सीकरण करता है, और उसके नाइट्रोजन को प्रोटॉनित करता है:
\ce{-NH3+} एक मेटा-निर्देशक, \omterm{def:b2:aromatic-substitution:directing}{निष्क्रियक समूह} है। ऐसीटैनिलाइड में नाइट्रोजन का एकाकी युग्म
कार्बोनिल के साथ साझा होता है: कम क्षारकीय और कम सक्रियक होने से न वह प्रोटॉनित होता है न ऑक्सीकृत,
और फिर भी ऑर्थो/पैरा निर्देशित करता है, त्रिविम कारणों से मुख्यतः पैरा।
\end{solution}

\begin{solution}{exo:b2:aromatic-substitution:9}
टॉलूईन का सल्फ़ोनीकरण कीजिए (मुख्यतः पैरा, भारी-भरकम समूह ऑर्थो से बचता है); ब्रोमीनन कीजिए: ब्रोमीन
मेथिल के ऑर्थो स्थान पर प्रवेश करता है (पैरा स्थिति अवरुद्ध है, और दोनों समूह उसी स्थिति के अनुकूल हैं);
तनु अम्ल में गर्म करके \ce{SO3H} समूह हटा दीजिए।
\end{solution}

\begin{solution}{exo:b2:aromatic-substitution:10}
दोनों समूह ऑर्थो/पैरा निर्देशित करते हैं; उनकी पैरा स्थितियाँ एक-दूसरे द्वारा घिरी हैं। मेथॉक्सी,
जो प्रबलतर सक्रियक है, जीतता है: \ce{-OCH3} के ऑर्थो स्थान पर नाइट्रीकरण (स्थिति 2), जिससे
4-मेथिल-2-नाइट्रोऐनिसोल मिलता है।
\end{solution}

\begin{solution}{exo:b2:aromatic-substitution:11}
हर नाइट्रो समूह वलय को प्रबल रूप से निष्क्रिय करता है: दूसरे नाइट्रीकरण को पहले से कठोर परिस्थितियाँ
चाहिए, और तीसरे को (ऐसे वलय पर जिस पर दो नाइट्रो समूह और केवल दुर्बल सक्रियक
मेथिल है) और भी कहीं अधिक कठोर।
\end{solution}

\begin{solution}{exo:b2:aromatic-substitution:12}
4-नाइट्रोबेन्ज़ोइक अम्ल: टॉलूईन का नाइट्रीकरण, पैरा समावयवी का पृथक्करण, मेथिल का
\ce{-COOH} में ऑक्सीकरण (गर्म परमैंगनेट)। 3-नाइट्रोबेन्ज़ोइक अम्ल: पहले टॉलूईन को बेन्ज़ोइक अम्ल में ऑक्सीकृत कीजिए, फिर
नाइट्रीकरण (\ce{-COOH} मेटा निर्देशित करता है)।
\end{solution}

\begin{solution}{pb:b2:aromatic-substitution:1}
\textbf{1.} \ce{HNO3 + 2H2SO4 -> NO2+ + H3O+ + 2HSO4-}।
\textbf{2.} \ce{O=N+=O}: रैखिक, \ce{CO2} के समइलेक्ट्रॉनी।
\textbf{3.} नाइट्रोजन \ce{CH3} के पैरा स्थित वलय कार्बन से आबंधित होता है, जो चतुष्फलकीय हो जाता है
(H और \ce{NO2} धारण करते हुए)।
\textbf{4.} धनात्मक आवेश आक्रांत कार्बन के ऑर्थो स्थित दोनों कार्बनों पर और उसके पैरा स्थित
कार्बन पर, जो मेथिल धारण करता है।
\textbf{5.} पहला, \omterm{def:b2:aromatic-substitution:seAr}{व्हीलैंड मध्यवर्ती} का बनना।
\textbf{6.} माध्यम का कोई क्षारक: \ce{HSO4-} या जल।
\textbf{7.} दो ऑर्थो, दो मेटा, एक पैरा।
\textbf{8.} 2 : 2 : 1, अर्थात् 40~\% ऑर्थो, 40~\% मेटा, 20~\% पैरा।
\textbf{9.} मेटा आक्रमण में धनात्मक आवेश कभी मेथिल धारण करने वाले कार्बन पर नहीं होता; किसी
संरचना को मेथिल से सहायता नहीं मिलती।
\textbf{10.} एक अनुनादी संरचना में आवेश \ce{CH3} धारण करने वाले कार्बन पर है: एक तृतीयक
कार्बधनायन, जो मेथिल के प्रेरणिक प्रभाव और \omterm{def:b2:fragment-orbitals:hyperconjugation}{अतिसंयुग्मन} से स्थायीकृत है।
\textbf{11.} ऑर्थो और पैरा अनुकूल हैं (96~\%, जबकि यादृच्छिक रूप से 60~\%); मेटा लगभग अनुपस्थित है।
\textbf{12.} प्रति स्थिति: ऑर्थो $58/2 = 29~\%$, पैरा 38~\%: पैरा अधिक अभिक्रियाशील स्थिति है,
क्योंकि ऑर्थो स्थितियाँ मेथिल द्वारा थोड़ी बाधित हैं।
\textbf{13.} तेज़: मेथिल वलय को सक्रिय करता है।
\textbf{14.} 4-नाइट्रोटॉलूईन (\qty{52}{\degreeCelsius} पर पिघलता है)।
\textbf{15.} अपरिष्कृत मिश्रण को ठंडा कीजिए: पैरा समावयवी क्रिस्टलित होता है और छानकर अलग किया जाता है;
द्रव में ऑर्थो और मेटा समावयवी रहते हैं।
\textbf{16.} सममित अणु क्रिस्टल में बेहतर संकुलित होता है: प्रति आयतन अधिक अन्योन्यक्रियाएँ,
अर्थात् उच्चतर गलनांक।
\textbf{17.} घटे हुए दाब पर \omterm{def:b2:liquid-vapour-diagrams:fractional-distillation}{प्रभाजी आसवन} से, या वर्णलेखन से; या
और ठंडा करके (मेटा समावयवी \qty{16}{\degreeCelsius} पर पिघलता है)।
\textbf{18.} समावयवियों के क्वथनांक निकट हैं (वही सूत्र, समान ध्रुवता)।
\textbf{19.} \ce{C7H8}: \qty{92.14}{g/mol}; \ce{C7H7NO2}: \qty{137.14}{g/mol}।
\textbf{20.} $100/92.14 = \qty{1.085}{mol}$ टॉलूईन; $0.90 \times 1.085 = \qty{0.977}{mol}$
नाइट्रोटॉलूईन।
\textbf{21.} कुल $0.977 \times 137.14 = \qty{134.0}{g}$: ऑर्थो \qty{77.7}{g}, मेटा
\qty{5.4}{g}, पैरा \qty{50.9}{g}।
\textbf{22.} दूसरा नाइट्रीकरण (डाइनाइट्रोटॉलूईन), और मेथिल समूह का ऑक्सीकरण।
\textbf{23.} गर्म क्षारीय परमैंगनेट से मेथिल को \ce{-COOH} में ऑक्सीकृत कीजिए, फिर अम्लीय बनाइए।
\textbf{24.} प्रयोग से 4-नाइट्रोटॉलूईन का $\boldsymbol{\approx \qty{51}{g}}$ मिलता है।
\end{solution}
